\textbf{Phần 1. Trắc nghiệm lựa chọn}
\setcounterpageref{ex}{0}

\begin{ex} %[3H9N50]
Chu vi của một hình là
\choice
{Phần mặt phẳng giới hạn bởi hình đó}
{Độ dài một cạnh}
{\True Tổng độ dài các cạnh bao quanh hình đó}
{Tích các cạnh của hình đó}
\loigiai{
Chu vi là tổng độ dài các cạnh bao quanh hình.\\
\textbf{Đáp án của câu hỏi là C}
}
\end{ex}

\begin{ex} %[3H9N50]
Muốn tính chu vi hình vuông ta lấy
\choice
{Độ dài một cạnh cộng với 4}
{Độ dài một cạnh nhân với chính nó}
{Độ dài một cạnh nhân với 2}
{\True Độ dài một cạnh nhân với 4}
\loigiai{
Chu vi hình vuông = Cạnh $\times 4$.\\
\textbf{Đáp án của câu hỏi là D}
}
\end{ex}

\begin{ex} %[3H9H50]
Chu vi hình tam giác có độ dài các cạnh $7\text{ cm}, 8\text{ cm}, 10\text{ cm}$ là
\choice
{\True $25\text{ cm}$}
{$23\text{ cm}$}
{$27\text{ cm}$}
{$30\text{ cm}$}
\loigiai{
$7 + 8 + 10 = 25\text{ cm}$.\\
\textbf{Đáp án của câu hỏi là A}
}
\end{ex}

\begin{ex} %[3H9H50]
Chu vi hình vuông có cạnh dài $12\text{ cm}$ là
\choice
{$36\text{ cm}$}
{$60\text{ cm}$}
{$24\text{ cm}$}
{\True $48\text{ cm}$}
\loigiai{
$12 \times 4 = 48\text{ cm}$.\\
\textbf{Đáp án của câu hỏi là D}
}
\end{ex}

\begin{ex} %[3H9H50]
Chu vi hình chữ nhật có chiều dài $15\text{ cm}$, chiều rộng $9\text{ cm}$ là
\choice
{$24\text{ cm}$}
{\True $48\text{ cm}$}
{$50\text{ cm}$}
{$60\text{ cm}$}
\loigiai{
$(15 + 9) \times 2 = 48\text{ cm}$.\\
\textbf{Đáp án của câu hỏi là B}
}
\end{ex}

\begin{ex} %[3H9H50]
Một hình tứ giác có độ dài 4 cạnh lần lượt là $5\text{ cm}, 6\text{ cm}, 7\text{ cm}, 8\text{ cm}$. Chu vi tứ giác đó là
\choice
{$24\text{ cm}$}
{$28\text{ cm}$}
{\True $26\text{ cm}$}
{$30\text{ cm}$}
\loigiai{
$5 + 6 + 7 + 8 = 26\text{ cm}$.\\
\textbf{Đáp án của câu hỏi là C}
}
\end{ex}

\begin{ex} %[3H9V50]
Một mảnh đất hình vuông có chu vi là $64\text{ m}$. Cạnh của mảnh đất đó dài là
\choice
{$14\text{ m}$}
{\True $16\text{ m}$}
{$18\text{ m}$}
{$20\text{ m}$}
\loigiai{
$64 : 4 = 16\text{ m}$.\\
\textbf{Đáp án của câu hỏi là B}
}
\end{ex}

\begin{ex} %[3H9V50]
Một hình chữ nhật có nửa chu vi là $25\text{ cm}$, chiều dài là $15\text{ cm}$. Chiều rộng hình chữ nhật là
\choice
{$8\text{ cm}$}
{$12\text{ cm}$}
{\True $10\text{ cm}$}
{$15\text{ cm}$}
\loigiai{
$25 - 15 = 10\text{ cm}$.\\
\textbf{Đáp án của câu hỏi là C}
}
\end{ex}

\begin{ex} %[3H9V50]
Một tam giác có chu vi $36\text{ cm}$ và 3 cạnh bằng nhau. Độ dài mỗi cạnh của tam giác là
\choice
{\True $12\text{ cm}$}
{$10\text{ cm}$}
{$14\text{ cm}$}
{$9\text{ cm}$}
\loigiai{
$36 : 3 = 12\text{ cm}$.\\
\textbf{Đáp án của câu hỏi là A}
}
\end{ex}

\begin{ex} %[3H9V50]
Một cái sân hình chữ nhật có chiều dài $20\text{ m}$, chiều rộng bằng nửa chiều dài. Chu vi của cái sân là
\choice
{$40\text{ m}$}
{$80\text{ m}$}
{$50\text{ m}$}
{\True $60\text{ m}$}
\loigiai{
Chiều rộng: $20 : 2 = 10\text{ m}$.\\
Chu vi: $(20 + 10) \times 2 = 60\text{ m}$.\\
\textbf{Đáp án của câu hỏi là D}
}
\end{ex}

\begin{ex} %[3H9V50]
Một khung ảnh hình vuông có cạnh dài $25\text{ cm}$. Người ta viền dây ruy-băng xung quanh khung ảnh. Chiều dài sợi dây ruy-băng cần dùng là
\choice
{\True $100\text{ cm}$ (hay $1\text{ m}$)}
{$80\text{ cm}$}
{$120\text{ cm}$}
{$90\text{ cm}$}
\loigiai{
$25 \times 4 = 100\text{ cm} = 1\text{ m}$.\\
\textbf{Đáp án của câu hỏi là A}
}
\end{ex}

\begin{ex} %[3H9C50]
Ghép 3 hình vuông cạnh $5\text{ cm}$ liền nhau thành một hình chữ nhật. Chu vi hình chữ nhật đó là
\choice
{$30\text{ cm}$}
{\True $40\text{ cm}$}
{$45\text{ cm}$}
{$60\text{ cm}$}
\loigiai{
Chiều dài: $5 \times 3 = 15\text{ cm}$, Chiều rộng: $5\text{ cm}$.\\
Chu vi: $(15 + 5) \times 2 = 40\text{ cm}$.\\
\textbf{Đáp án của câu hỏi là B}
}
\end{ex}

\textbf{Phần 2. Trắc nghiệm đúng sai}
\setcounterpageref{ex}{0}

\begin{ex} %[3H9H50]
Xét tính đúng sai về công thức tính chu vi:
\choiceTF
{\True Chu vi hình chữ nhật = $(\text{Chiều dài} + \text{Chiều rộng}) \times 2$.}
{\True Chu vi hình vuông = $\text{Độ dài cạnh} \times 4$.}
{Hình chữ nhật có chiều dài $8\text{ cm}$, rộng $4\text{ cm}$ có chu vi là $24\text{ cm}$.}
{\True Chu vi hình chữ nhật đó là $(8 + 4) \times 2 = 24\text{ cm}$.}
\loigiai{
\begin{itemize}
    \item Công thức chu vi chữ nhật: Đúng.
    \item Công thức chu vi hình vuông: Đúng.
    \item Tính toán: $(8 + 4) \times 2 = 24\text{ cm}$. Cả hai ý cuối đều Đúng.
\end{itemize}
\begin{itemize}
    \item a) \textbf{Đúng}.
    \item b) \textbf{Đúng}.
    \item c) \textbf{Đúng}.
    \item d) \textbf{Đúng}.
\end{itemize}
}
\end{ex}

\begin{ex} %[3H9V50]
Cho hình chữ nhật có chu vi $36\text{ cm}$, chiều rộng $6\text{ cm}$. Xét tính đúng sai:
\choiceTF
{\True Nửa chu vi hình chữ nhật là $18\text{ cm}$.}
{\True Chiều dài hình chữ nhật là $12\text{ cm}$.}
{Chiều dài gấp 3 lần chiều rộng.}
{\True Chiều dài gấp 2 lần chiều rộng ($12 : 6 = 2$).}
\loigiai{
\begin{itemize}
    \item Nửa chu vi: $36 : 2 = 18\text{ cm}$. Đúng.
    \item Dài: $18 - 6 = 12\text{ cm}$. Đúng.
    \item $12 : 6 = 2$ lần $\neq 3$ lần. Ý 3 Sai, Ý 4 Đúng.
\end{itemize}
\begin{itemize}
    \item a) \textbf{Đúng}.
    \item b) \textbf{Đúng}.
    \item c) \textbf{Sai}.
    \item d) \textbf{Đúng}.
\end{itemize}
}
\end{ex}

\textbf{Phần 3. Trả lời ngắn}
\setcounterpageref{ex}{0}

\begin{ex} %[3H9H50]
Tính chu vi hình vuông có cạnh dài $14\text{ cm}$.
\shortans{56}
\loigiai{
\textbf{Đáp số:} $56$\\
$14 \times 4 = 56\text{ cm}$.
}
\end{ex}

\begin{ex} %[3H9H50]
Tính chu vi hình chữ nhật có chiều dài $22\text{ cm}$, chiều rộng $8\text{ cm}$.
\shortans{60}
\loigiai{
\textbf{Đáp số:} $60$\\
$(22 + 8) \times 2 = 60\text{ cm}$.
}
\end{ex}

\begin{ex} %[3H9V50]
Một hình vuông có chu vi $48\text{ dm}$. Tính cạnh của hình vuông theo đề-xi-mét.
\shortans{12}
\loigiai{
\textbf{Đáp số:} $12$\\
$48 : 4 = 12\text{ dm}$.
}
\end{ex}

\begin{ex} %[3H9V50]
Tính chu vi hình tứ giác có 4 cạnh đều bằng $9\text{ cm}$.
\shortans{36}
\loigiai{
\textbf{Đáp số:} $36$\\
$9 \times 4 = 36\text{ cm}$.
}
\end{ex}

\begin{ex} %[3H9V50]
Một hình chữ nhật có chu vi $50\text{ cm}$, chiều dài là $18\text{ cm}$. Tính chiều rộng hình chữ nhật theo xăng-ti-mét.
\shortans{7}
\loigiai{
\textbf{Đáp số:} $7$\\
Nửa chu vi: $25\text{ cm}$. Chiều rộng: $25 - 18 = 7\text{ cm}$.
}
\end{ex}

\begin{ex} %[3H9C50]
Uốn một sợi dây đồng vừa đủ thành hình vuông cạnh $15\text{ cm}$. Nếu uốn sợi dây đó thành hình chữ nhật có chiều dài $18\text{ cm}$ thì chiều rộng là bao nhiêu xăng-ti-mét?
\shortans{12}
\loigiai{
\textbf{Đáp số:} $12$\\
Độ dài sợi dây: $15 \times 4 = 60\text{ cm}$.\\
Nửa chu vi: $60 : 2 = 30\text{ cm}$. Chiều rộng: $30 - 18 = 12\text{ cm}$.
}
\end{ex}
